\SourceLeaf{72}nite parva figurae fluidi, pro qua volumen $s$ invariatum manet, i. e. pro qua
\[
\delta s=\int\mathrm{d}U.\delta e.\cos(4,5)
\]
evanescit, facile perspicitur, figuram superficiei $U$ in statu aequilibrii talem esse debere, ut in omnibus eius punctis elementum variationis $\delta W$ hoc
\[
\mathrm{d}U.\delta e.\cos(4,5).\bigl[z+\alpha\alpha(\tfrac1R+\tfrac1{R'})\bigr]
\]
proportionale sit elemento variationis $\delta s$, puta quantitati $\mathrm{d}U.\delta e.\cos(4,5)$, sive quod idem est, ut fiat
\[
z+\alpha\alpha(\tfrac1R+\tfrac1{R'})=\mathrm{Const.}
\]
Manifesto enim, si haec proportionalitas locum non haberet, valor ipsius $W$ decrementi capax foret per idoneam mutationem figurae superficiei $U$, limite $P$ adeo invariato manente. Ceterum aequatio illa pro \textit{tota} superficie $U$ valet, etiamsi haec e pluribus partibus separatis constet, dummodo fluidum ipsum cohaereat.

Aequatio ista constituit theorema fundamentale primum in theoria aequilibrii fluidorum, quod iam ab ill. \textsc{Laplace} erutum est, sed per methodum a nostra plane diversam.

Si planum, pro quo $z$ quantitati aequationis constanti aequalis est. et quod planum horizontale normale (plan de niveau) vocare possumus, loco eius, a quo coordinatae $z$ numeratae erant, adoptamus, erit
\[
z=-\alpha\alpha(\tfrac1R+\tfrac1{R'})
\]
unde protinus demanant corollaria sequentia.

I. Si planum normale superficiem liberam $U$ ullibi secat, in quovis sectionis puncto superficies necessario concavo-convexa erit, atque radius maximus convexitatis radio maximo concavitatis aequalis.

II. Supra planum normale superficies vel concavo-concava erit, vel, sicubi fuerit concavo-convexa, curvatura concava convexam superabit.

III. Infra planum normale superficies vel erit convexo-convexa, vel sicubi fuerit concavo-convexa, curvatura convexa concavam superabit.

IV. Superficies libera $U$ nequit habere partem finitam planam nisi horizontalem et cum plano normali coincidentem.

\SourceLeaf{73}\Article{29}
Aequatione, quam modo stabilivimus, subsistente, variatio valoris ipsius $W$ reducitur ad
\[
\delta W=-\int\mathrm{d}P.\delta e.\cos(5,8)(\alpha\alpha\cos i-\alpha\alpha+2\beta\beta)
\]
unde introducendo angulum $A$ talem ut sit
\[
\cos A=\frac{\alpha\alpha-2\beta\beta}{\alpha\alpha}\qquad\text{sive}\qquad\sin\tfrac12 A=\frac\beta\alpha
\]
habemus
\[
\delta W=\alpha\alpha\int\mathrm{d}P.\delta e.\cos(5,8).(\cos A-\cos i)
\]
integratione per totam lineam $P$ extensa. Memores esse debemus, factorem $\cos(5,8)$ aequalem esse ipsi $\sin(5,6)$, signo positivo vel negativo affecto, prout fluidum in motu suo virtuali apud elementum $\mathrm{d}P$ vel ultra limitem $P$ redundare, vel citra recedere concipitur. Hinc facile concludimus, in statu aequilibrii, generaliter loquendo, ubique esse debere $i=A$. Si enim in aliqua parte lineae $P$ esset $i<A$, motus virtualis primi generis in hac parte, manente parte reliqua limitis $P$ invariata, manifesto ipsi $W$ variationem negativam induceret, et perinde negativa variatio ipsius $W$ prodiret per motum virtualem fluidi secundi generis, si in ulla parte lineae $P$ esset $i>A$: utraque igitur suppositio conditioni minimi in aequilibrio adversatur.

Hoc est theorema fundamentale secundum, quod etiam investigationibus ill. \textsc{Laplace} intertextum, sed e principio virium molecularium haud demonstratum videmus.

\Article{30}
Theorema art. praec. modificatione quadam eget in casu singulari, quem silentio praeterire non licet. Tacite scilicet supposuimus, superficiem vasis iuxta totum limitem $P$ curvatura continua gaudere, ita ut in quovis huius limitis puncto \textit{unicum} planum superficiem vasis tangens exstet. Si continuitas curvaturae in aliquo puncto singulari lineae $P$ interrumpitur, sive cuspis ibi adsit, sive acies lineam $P$ traiiciens, facile perspicietur, conclusionem nostram hinc non immutari; sed aliter res se habet. si continuitas curvaturae interrupta est in parte finita lineae $P$. i. e. si superficies vasis per partem finitam lineae $P$ (vel adeo per totam hanc lineam) aciem offert. Tunc scilicet in quovis talis partis puncto bina plana
\SourceLeaf{74}superficiem vasis tangentia aderunt, quorum alterum refertur ad partem liberam superficiei vasis, alterum ad partem $T$. Retinendo itaque characterem $i$ pro angulo inter planum prius atque planum tangens superficiem $U$, denotandoque per $k$ angulum inter hoc planum et planum posterius, haud amplius erit $i+k=180^0$, sed maior minorve, prout acies est convexa vel concava. Et dum elementum variationis $\delta W$, pro motu virtuali fluidi ultra limitem $P$ redundantis, etiamnum exprimitur per
\[
\alpha\alpha\mathrm{d}P.\delta e.\sin(5,6).(\cos A-\cos i)
\]
elementum illius variationis pro motu virtuali fluidi citra limitem $P$ recedentis iam erit
\[
-\alpha\alpha\mathrm{d}P.\delta e.\sin(5,6).(\cos A+\cos k)
\]
Ne igitur valor ipsius $W$ capax sit variationis negativae, requiritur, ut neque valor ipsius $\cos A-\cos i$ sit negativus, neque valor ipsius $\cos A+\cos k$ positivus, i. e. esse debet
\[
\begin{aligned}
&\text{vel}\quad i=A,\qquad\text{vel}\quad i>A\\
&\text{atque vel}\quad k=180^0-A,\qquad\text{vel}\quad k>180^0-A
\end{aligned}
\]
In statu aequilibrii itaque esse nequit $i+k<180^0$, sive, quod idem est, \textit{in statu aequilibrii limes superficiei fluidi liberae $U$ esse nequit, per extensionem finitam, in acie concava superficiei vasis.} Contra, quoties pars illius limitis coincidit cum acie convexa, ad aequilibrium requiritur et sufficit, ut angulus inter plana fluidum et vas tangentia sit inter limites $A$ et $A+a$ (incl.), extra fluidum, sive inter $180^0-A$ et $180^0-A+a$, intra fluidum mensuratus, si angulum inter duo plana superficiem vasis utrimque ab acie tangentia in quovis puncto indefinite per $180^0-a$ denotamus, quatenus hic angulus a plaga vasis mensuratur.

\Article{31}
Constantes $\alpha\alpha$, $\beta\beta$, quarum ratio angulum $A$ determinat, a functionibus $f,F$ pendent, et quodammodo tamquam mensurae intensitatis virium molecularium, quas particulae fluidi et vasis exercent, considerari possunt. Si functiones istae ita comparatae sunt, ut $fx,Fx$ sint in ratione determinata a distantia $x$ independente, puta ut $n$ ad $N$, manifesto statuere possumus
\SourceLeaf{75}$\alpha\alpha:\beta\beta=cn:CN$, i. e. constantes $\alpha\alpha$, $\beta\beta$ erunt proportionales attractionibus, quas in eadem distantia exercent duae moleculae quoad volumen aequales, altera fluidi altera vasis. Iam quum angulus $A$ fiat acutus, recto aequalis, obtusus, duobus rectis aequalis, prout $\beta\beta<\frac12\alpha\alpha$, $\beta\beta=\frac12\alpha\alpha$, $\beta\beta>\frac12\alpha\alpha$ sed $<\alpha\alpha$, $\beta\beta=\alpha\alpha$: in sensu istius suppositionis (quae si gratuita est, tamen verisimilitudini non repugnat) dicere oportet, casum primum locum habere, quoties attractio partium fluidi mutua maior sit quam duplum attractionis partium vasis in fluidum; secundum, quoties prior attractio sit duplum posterioris; tertium, quoties prior maior quidem sit posteriori, sed minor eius duplo; denique quartum, quoties ambae attractiones sint aequales. Exemplum casus primi exhibet argentum vivum in vasibus vitreis.

\Article{32}
At quantus est valor anguli $A$ in casu eo, ubi attractio vasis maior est quam attractio partium fluidi mutua? Valor imaginarius, quem pro $\beta\beta>\alpha\alpha$ formula $\sin\frac12 A=\frac\beta\alpha$ angulo $A$ assignat, iam testatur, suppositionem aliquam in tali casu non admissibilem subesse. Revera quoties $\beta\beta>\alpha\alpha$, suppositio \textit{limitationis} superficiei $T$ cum conditione minimi respectu functionis $W$ consistere nequit. Ubicunque enim limitem posueris, patet, si ultra hunc limitem cutem fluidi tenuissimam expansam concipias, ita ut $T$ capiat augmentum $T'$, et proin $U$ augmentum huic proxime aequale, valorem functionis $W$ assumere mutationem sensibiliter aequalem quantitati negativae $-(2\beta\beta-2\alpha\alpha)T'$; quinadeo valorem ipsius $W$ tamdiu ulterioris diminutionis capacem manere, donec $T'$ totam superficiem vasis reliquam occupaverit. Valor mutationis $-(2\beta\beta-2\alpha\alpha)T'$ eo magis exactus erit, quo minor crassities accipiatur, et quatenus tantummodo de valore expressionis $W$ agitur, nihil impedit, quominus crassities usque ad evanescendum diminui concipiatur. Attamen cutis crassitiei evanescentis (probe distinguendae ab insensibili) nihil esset nisi fictio mathematica, figuraque spatii $s$ tali fictioni accommodata revera non differret ab ea, pro qua $W$ in casu $\beta\beta=\alpha\alpha$ valorem minimum acquirit.

Sed paullo aliter res se habet in problemate nostro physico, ubi talis cutis accessoria necessario gaudere debet certa crassitie, utut insensibili, quo aequilibrium consistere possit. Quoties talis pars adest, expressio $W$, uti in art. 18 docuimus, incompleta est, et denotata ea parte vasis, quam cutis tegit, per $T'$,
\SourceLeaf{76}huiusque crassitie in quovis puncto indefinite per $\rho$, expressioni $\Omega$ adhuc adiiciendi erunt termini
\[
\pi cc\int\theta'\rho.\mathrm{d}T'-\pi cC\int\Theta'\rho.\mathrm{d}T'
\]
adeoque valori ipsius $W$ hi
\[
\begin{gathered}
\frac{\pi C}g\int\Theta'\rho.\mathrm{d}T'-\frac{\pi c}g\int\theta'\rho.\mathrm{d}T'\\
=\int\mathrm{d}T'.\left(\frac{2\beta\beta}{\Theta0}.\Theta'\rho-\frac{2\alpha\alpha}{\theta0}.\theta'\rho\right)
\end{gathered}
\]
Quocirca quum valor ipsius $W$, per accessionem istius cutis, iam acceperit mutationem $(2\beta\beta-2\alpha\alpha)T'$, mutatio tota, ei valori ipsius $W$, qui omittendo cutem locum habet, adiicienda, erit
\[
-2\int\mathrm{d}T'.\left[\beta\beta\left(1-\frac{\Theta'\rho}{\Theta0}\right)-\alpha\alpha\left(1-\frac{\theta'\rho}{\theta0}\right)\right]
\]
Haec mutatio propter $\theta'0=\theta0$, $\Theta'0=\Theta0$, nulla esset pro crassitie evanescente; at quum $\theta'\rho$, $\Theta'\rho$, crescente crassitie $\rho$, citissime decrescant, et iam pro valore insensibili ipsius $\rho$ insensibiles evadant, mutatio ista citissime versus valorem $-(2\beta\beta-2\alpha\alpha)T'$ converget, atque pro statu aequilibrii fluidi, ne valor expressionis $W$ correctae capax sit ulterioris diminutionis sensibilis, sensibiliter eidem aequalis esse debebit. Ceterum investigatio completa legis, quam crassities $\rho$ sequi debet, profundiores evolutiones requireret, quibus tamen hic non immoramur, quum absque cognitione functionum $f,F$, a quibus functiones $\theta',\Theta'$ pendent, nec non propter rationes in art. 34 indicandas, nimis otiosae videri possent. Ad investigationem partis substantialis fluidi, i. e. eius, cuius dimensiones omnes sensibiles sunt, sufficit, pro casu nostro, ubi $\beta\beta>\alpha\alpha$, vas in vicinia limitis partis substantialis \textit{madefactum} concipere, i. e. cute fluida obductam, cuius crassities insensibilis quidem sit, attamen tanta, ut $\theta'\rho$, $\Theta'\rho$ negligi possint. Hoc pacto functio, quae in statu aequilibrii minimum esse debet, erit
\[
\int z\mathrm{d}s-2(\beta\beta-\alpha\alpha)(T+T')-\alpha\alpha T+\alpha\alpha U
\]
ubi $T,U$ ad solam partem substantialem fluidi referri supponuntur. Patet itaque, variationem huius functionis e mutatione virtuali figurae partis substantialis fluidi oriundam (qualis mutatio aggregatum $T+T'$ non afficit) convenire cum variatione expressionis
